\section{规模不是一个数字，而是多个不可混用的队列}

上一节确定了响应链，本节进一步解释为什么“规模增长”会压垮组织。公开统计常同时出现 reports、files、accounts、incidents 和 victims；它们的单位、重复率与年度口径不同。工程团队若只把最大的数字写进 dashboard，就会错误估计存储、审核、调查和模型训练负担。

\subsection{Reports、files、duplicates 与 novel candidates}

一个 report 可能包含多个文件，同一文件也可能被不同平台或账户多次报告；而 investigators 真正需要优先处理的，往往是可能对应新受害、持续风险或此前未知材料的部分。\term{revictimization}（再次伤害）指受害材料被持续复制、观看或传播给受害者带来的重复伤害，因此即使是“已知文件”的重复出现也值得快速阻断，但不应与“新案件数量”混为一谈。

\lecturefigure{02-scale-funnel.png}{报告量、文件量、重复内容、未知候选和稀缺复核能力是不同度量。}{本地字幕 00:00--05:50；NCMEC 年度口径说明；概念重绘。}

\begin{knowledgebox}{读图：先问分母和去重规则}
读任何规模数据时先写出单位：report 是一次提交，file 是附件，candidate 是算法筛出的待复核对象，case 是机构按规则聚合后的调查单元。还要问是否按 cryptographic hash、perceptual hash、account、事件或 victim 去重。NCMEC 在不同年度调整过统计与报告方法，因此 2023、2024、2025 页面只能在各自定义内解释，不能把曲线转折直接归因于真实危害同比变化。
\end{knowledgebox}

\begin{importantbox}{容量规划要同时看四种率}
设单位时间到达的报告数为 $\lambda_r$，每份报告平均文件数为 $m_f$，已知内容比例为 $p_k$，需要人工复核的未知候选比例为 $p_u$，则粗略审核到达率为
\[
\lambda_h \approx \lambda_r m_f p_u.
\]
其中 $\lambda_h$ 不是受害者数量，也不是法律案件数量；它只是审核队列的负载估计。已知匹配仍会触发处置与报告，因此不能从人工队列中消失就误认为没有成本。
\end{importantbox}

\teachervoice{课堂反复区分平台的“检测并报告”与调查人员的“识别受害者并追查责任人”。讲者强调，后端机构最不需要的是另一个没有排序、没有上下文、包含大量重复项的文件堆。}

\subsection{本章小结}

规模问题的本质是异质队列：重复传播要高速阻断，未知内容要谨慎复核，高优先级线索要携带上下文，调查人员则需要可行动的 case package。单一总量无法指导系统设计。

